% from main.tex
%%%%%%%%%%  scan idx 567  |  folio 554  %%%%%%%%%%
\textbf{Proposition 6.2.2.} Pour toute section positive \(n\) de \(\mathbb Z\), le faisceau étale en groupes \(\Der(\Gl(n))\) des commutateurs de \(\Gl(n))\) est représentable par
\(\operatorname{Sl}(n)=\Ker(\det:\Gl(n)\to\Gm)\).

\textbf{Corollaire 6.2.3.} Toute représentation
\begin{equation}\tag{6.2.4}
 \alpha:G=\Gl(n_1)\times\cdots\times\Gl(n_r)\longrightarrow\Gm
\end{equation}
est de la forme
\begin{equation}\tag{6.2.5}
 \alpha(Y)(s)=(\det(s_1))^{m_1}\cdots(\det(s_r))^{m_r},
\end{equation}
où \(Y\in\operatorname{Ob}(\Sch/S)\), \(s=(s_1,\ldots,s_r)\in G(Y)\) et \(m_1,\ldots,m_r\) sont des sections de \(\mathbb Z\) uniquement déterminés par \(\alpha\).

\textbf{Corollaire 6.2.6.} Avec les notations de 6.1. on a
\begin{equation}\tag{6.2.6}
 \det\bigl(T_\alpha(E_1,\ldots,E_r)\bigr)
 =
 \det(E_1)^{\otimes m_1}\otimes\cdots\otimes
 \det(E_r)^{\otimes m_r},
\end{equation}
où \(m_1,\ldots,m_r\) sont des sections de \(\mathbb Z\) uniquement déterminées par \(\alpha\).

\textbf{Remarque 6.2.7.} On peut déterminer les exposants \(m_1,\ldots,m_r\) dans (6.2.5) et (6.2.6) en substituant successivement dans (6.2.5)
\(s_i=\lambda.\ \operatorname{id}.(\lambda\in\Gm(Y))\), \(s_j=\operatorname{id}_{\Gl(n_j)}(Y)\), pour \(j\ne i\) \((i,j=1,\ldots,r)\), d'où
\[
 \det(s_i)=\lambda^{n_i}\quad\text{et}\quad
 \alpha(Y)(s)=\lambda^{n_im_i}.
\]
Par exemple on déduit ainsi de (6.2.6)
\begin{align}
 \det(E_i\otimes E_2)&=(\det(E_1))^{\otimes n_2}\otimes(\det(E_2))^{\otimes n_1}, \tag{6.2.8}\\
 \det(E_1\otimes E_2)&=(\det(E_1))\otimes(\det(E_2)) \tag{6.2.9}
\end{align}
et
\begin{equation}\tag{6.2.10}
 \det\bigl(\bigwedge\limits^m E_1\bigr)=(\det(E_1))^{\binom{n_1-1}{m-1}}
\end{equation}
etc.

\subsection*{6.3}
